\section*{14. 天际轮廓问题}

给定n座建筑物，每个建筑物用三元组(Li,Ri,Hi)表示（Li为左下x坐标，Ri为右下x坐标，Hi为高度），设计O(nlogn)算法求出这n座建筑物的天际轮廓。

\begin{itemize}
    \item \textbf{天际轮廓表示}：用高度变化的关键点序列表示，如单个建筑物的轮廓为[(Li,0),(Li,Hi),(Ri,Hi),(Ri,0)]
    \item \textbf{分治分解}：将建筑物集均分为左右两个子集，递归求得左右
    \item \textbf{双指针合并}：用双指针i、j分别遍历左右天际线，按x坐标升序合并，同时维护当前左右高度hL、hR，取max(hL,hR)作为合并后当前高度
    \item \textbf{关键点记录}：当当前max(hL,hR)与上次记录的高度不同时，记录新的关键点(x, max(hL,hR))
    \item \textbf{时间复杂度}$T(n)=2T(n/2)+O(n)$，得$O(n\log n)$
\end{itemize}

\begin{lstlisting}[language=C++, style=cppstyle]
struct Point {
    int x, height;
    bool isStart;
};
vector<Point> mergeSkylines(vector<Point>& left, vector<Point>& right) {
    vector<Point> result;
    int i = 0, j = 0;
    int hL = 0, hR = 0, currH = 0;
    while (i < left.size() && j < right.size()) {
        int x;
        if (left[i].x < right[j].x ||
            (left[i].x == right[j].x && left[i].height >= right[j].height)) {
            hL = left[i].height;
            x = left[i].x;
            i++;
        } else {
            hR = right[j].height;
            x = right[j].x;
            j++;}}
        int newH = max(hL, hR);
        if (result.empty() || newH != currH) {
            result.push_back({x, newH});
            currH = newH;}}
    while (i < left.size()) {
        if (left[i].height != currH) {
            result.push_back(left[i]);
            currH = left[i].height;}}
        i++;}
    while (j < right.size()) {
        if (right[j].height != currH) {
            result.push_back(right[j]);
            currH = right[j].height;}}
        j++;}
    return result;}
vector<Point> getSkyline(vector<vector<int>>& buildings, int low, int high) {
    if (low == high) {
        return {{buildings[low][0], buildings[low][2]},
                {buildings[low][1], 0}};}}
    int mid = (low + high) / 2;
    auto left = getSkyline(buildings, low, mid);
    auto right = getSkyline(buildings, mid + 1, high);
    return mergeSkylines(left, right);}}
\end{lstlisting}

\begin{enumerate}
    \item 如果只有一个建筑，天际线就是它的4个顶点
    \item 多个建筑时，从中间分成左右两半，递归求各自的天际线
    \item 用双指针遍历两条天际线，x坐标小的先处理
    \item 维护当前左右天际线的高度，取最大值为合并后当前高度
    \item 只要当前最大高度和上次记录不一样，就记录新的关键点
    \item 重复直到处理完所有关键点，最终得到的就是完整天际线
\end{enumerate}
